\section{Introduction}
\label{sec:intro}

Short-term photovoltaic (PV) power forecasting is an active research
area, and reported accuracy improvements for hybrid deep-learning
architectures are substantial. We built an open, leakage-controlled
benchmark harness for this task on a public dataset: five
architectures, evaluated across three module technologies, three
forecast horizons, two feature regimes and five random seeds, with 900
machine-readable run records that reproduce every table and figure in
this paper.

To characterise how such gains are currently reported, we coded 27
papers on hybrid PV power forecasting, published between 2022 and
2026, on eight dimensions of evaluation protocol, each judgement
supported by a verbatim quotation from the source. Two counts
summarise the picture: 26 of 27 report no skill score against any
reference forecast, and all 27 report point estimates with no measure
of run-to-run variance. Reported numbers in this literature are
therefore difficult to compare with one another, or with a naive
forecast, without knowing what they were measured against.

This paper measures what that comparison is worth. Holding the model,
data and training procedure fixed, we vary only the reference forecast
and the treatment of night hours. On one array at a six-hour horizon,
on the held-out test year, the same forecasts score 0.701 against
smart persistence and 0.232 against the optimal convex combination of
climatology and persistence recommended by Yang et
al.~\cite{yang2020verification} - 67 percent of the apparent skill
relative to persistence disappears when the same forecasts are
evaluated against the convex reference (70 percent on the validation
split, an upper bound, since the reference's weight is fitted there).
More consequentially, skill against persistence rises monotonically
with horizon, the pattern usually read as a model's advantage growing
at longer lead times, while skill against the recommended reference is
non-monotonic and peaks at three hours: the trend is a property of the
reference, not the model. Including night hours deflates normalised
RMSE by approximately 34 percent, a factor that is closed-form in the
daylight fraction, while the skill score is nearly unchanged. Among
the three base architectures evaluated under an identical protocol,
the largest difference we measure on the held-out test year is 0.028
in skill score - smaller than either protocol effect.

\subsection{Contributions}
\label{sec:intro-contributions}

1. An open, leakage-controlled benchmark harness for short-term PV
power forecasting, with 900 machine-readable run records reproducing
every reported number.

2. A measurement of two evaluation-protocol effects - reference
forecast and night-hour inclusion - held against a fixed model and
data, and shown to exceed the largest architectural difference we
measure.

3. A coded survey of 27 papers characterising current reporting
practice on reference forecasts and run-to-run variance.

This paper is organised around three questions. RQ1: how much does the
choice of reference forecast change the magnitude and shape of
reported skill (Section~\ref{sec:rq1-reference})? RQ2: how much does
including night hours change reported error and skill
(Section~\ref{sec:rq1-night})? RQ3: how large are these protocol
effects relative to the largest difference we measure between
architectures (Section~\ref{sec:rq2-architecture})?
